\section{Bài tập 3: Circuit Analysis with Diode}

Cho mạch điện hình cầu chứa diode và các điện trở như hình bên dưới. Mạch gồm nguồn điện áp $V$, các điện trở $R_1 = 2\text{ k}\Omega$, $R_2 = 8\text{ k}\Omega$, $R_3 = 4\text{ k}\Omega$, $R_4 = 6\text{ k}\Omega$ và diode $D_3$ đặt ở nhánh cầu giữa nối từ nút $1$ (anode) sang nút $2$ (cathode). 

Hãy phân tích mạch bằng mô hình diode thực tế với điện áp rơi khi phân cực thuận là $V_D = 0.7\text{ V}$. Thiết lập hệ phương trình để xác định dòng điện qua tất cả các điện trở ($I_{R1}, I_{R2}, I_{R3}, I_{R4}$) cùng điện thế nút $V_1, V_2$ trong hai trường hợp điện áp nguồn:
\begin{itemize}[leftmargin=1.5em]
  \item Trường hợp 1: $V = 8\text{ V}$.
  \item Trường hợp 2: $V = 12\text{ V}$.
\end{itemize}
Sau đó, tiến hành mô phỏng PSpice ở chế độ Bias Point để kiểm chứng kết quả.

\given{ex3_given}{0.65\textwidth}{Mạch cầu phân tích với Diode (Hình 2.16)}

\subsection{Tính toán lý thuyết}
Gọi $V_1$ là điện thế tại nút anode của diode (giữa $R_1$ và $R_2$) và $V_2$ là điện thế tại nút cathode của diode (giữa $R_3$ và $R_4$). Chọn nút cực âm của nguồn $V$ làm mass tham chiếu ($0\text{ V}$).

\noindent Khi ngắt diode:
\begin{align*}
  V_{1,\text{open}} &= V \cdot \frac{R_2}{R_1 + R_2} = V \cdot \frac{8}{2 + 8} = 0.8\,V \\
  V_{2,\text{open}} &= V \cdot \frac{R_4}{R_3 + R_4} = V \cdot \frac{6}{4 + 6} = 0.6\,V \\
  V_{1,\text{open}} - V_{2,\text{open}} &= 0.8\,V - 0.6\,V = 0.2\,V
\end{align*}
Với cả hai giá trị nguồn $V = 8\text{ V}$ ($0.2 \times 8 = 1.6\text{ V} > 0.7\text{ V}$) và $V = 12\text{ V}$ ($0.2 \times 12 = 2.4\text{ V} > 0.7\text{ V}$), hiệu điện thế thuận đều lớn hơn điện áp ngưỡng $0.7\text{ V}$. Do đó diode $D_3$ phân cực thuận.

\noindent Theo mô hình thực tế, hiệu điện thế giữa hai nút thỏa mãn:
\begin{equation}
  V_1 - V_2 = V_D = 0.7\text{ V} \iff V_1 = V_2 + 0.7 \label{eq:ex3_diode}
\end{equation}

\noindent Áp dụng định luật KCL tại nút $1$:
\begin{equation}
  I_{R1} = I_{D3} + I_{R2} \implies \frac{V - V_1}{R_1} = I_{D3} + \frac{V_1}{R_2} \label{eq:ex3_kcl1}
\end{equation}

\noindent Áp dụng định luật KCL tại nút $2$:
\begin{equation}
  I_{R3} + I_{D3} = I_{R4} \implies I_{D3} = \frac{V_2}{R_4} - \frac{V - V_2}{R_3} \label{eq:ex3_kcl2}
\end{equation}

\noindent Cộng vế theo vế hai phương trình (\ref{eq:ex3_kcl1}) và (\ref{eq:ex3_kcl2}) để triệt tiêu dòng $I_{D3}$:
\begin{equation}
  \frac{V - V_1}{R_1} + \frac{V - V_2}{R_3} = \frac{V_1}{R_2} + \frac{V_2}{R_4}
\end{equation}
\begin{equation}
  V\left(\frac{1}{R_1} + \frac{1}{R_3}\right) = V_1\left(\frac{1}{R_1} + \frac{1}{R_2}\right) + V_2\left(\frac{1}{R_3} + \frac{1}{R_4}\right)
\end{equation}

\noindent Thay $V_1 = V_2 + 0.7$ vào phương trình trên:
\begin{equation}
  V\left(\frac{1}{2000} + \frac{1}{4000}\right) = (V_2 + 0.7)\left(\frac{1}{2000} + \frac{1}{8000}\right) + V_2\left(\frac{1}{4000} + \frac{1}{6000}\right)
\end{equation}
\begin{equation}
  \frac{3V}{4000} = \frac{5}{8000}(V_2 + 0.7) + \frac{5}{12000}V_2 \iff 0.75 \times 10^{-3} V = 0.625 \times 10^{-3}(V_2 + 0.7) + 0.4167 \times 10^{-3} V_2
\end{equation}
\begin{equation}
  1.0417 \times 10^{-3} V_2 = 0.75 \times 10^{-3} V - 0.4375 \times 10^{-3}
\end{equation}

\subsubsection*{1. Với $V = 8\text{ V}$:}
\begin{align*}
  V_2 &= \frac{0.75 \times 10^{-3} \times 8 - 0.4375 \times 10^{-3}}{1.0417 \times 10^{-3}} = \frac{6 - 0.4375}{1.04167} = 5.3400\text{ V} \\
  V_1 &= V_2 + 0.7 = 5.3400 + 0.7 = 6.0400\text{ V} \\
  I_{R1} &= \frac{8 - 6.04}{2\text{ k}\Omega} = \frac{1.96}{2000} = 0.9800\text{ mA} \\
  I_{R2} &= \frac{6.04}{8\text{ k}\Omega} = 0.7550\text{ mA} \\
  I_{R3} &= \frac{8 - 5.34}{4\text{ k}\Omega} = \frac{2.66}{4000} = 0.6650\text{ mA} \\
  I_{R4} &= \frac{5.34}{6\text{ k}\Omega} = 0.8900\text{ mA} \\
  I_{D3} &= I_{R1} - I_{R2} = 0.9800 - 0.7550 = 0.2250\text{ mA}
\end{align*}

\subsubsection*{2. Với $V = 12\text{ V}$:}
\begin{align*}
  V_2 &= \frac{0.75 \times 10^{-3} \times 12 - 0.4375 \times 10^{-3}}{1.0417 \times 10^{-3}} = \frac{9 - 0.4375}{1.04167} = 8.2200\text{ V} \\
  V_1 &= V_2 + 0.7 = 8.2200 + 0.7 = 8.9200\text{ V} \\
  I_{R1} &= \frac{12 - 8.92}{2\text{ k}\Omega} = \frac{3.08}{2000} = 1.5400\text{ mA} \\
  I_{R2} &= \frac{8.92}{8\text{ k}\Omega} = 1.1150\text{ mA} \\
  I_{R3} &= \frac{12 - 8.22}{4\text{ k}\Omega} = \frac{3.78}{4000} = 0.9450\text{ mA} \\
  I_{R4} &= \frac{8.22}{6\text{ k}\Omega} = 1.3700\text{ mA} \\
  I_{D3} &= I_{R1} - I_{R2} = 1.5400 - 1.1150 = 0.4250\text{ mA}
\end{align*}

\begin{dapsobox}
  \begin{itemize}[leftmargin=1.5em]
    \item Khi $V = 8\text{ V}$: $V_1 = 6.04\text{ V}$, $V_2 = 5.34\text{ V}$, $I_{R1} = 0.98\text{ mA}$, $I_{R2} = 0.755\text{ mA}$, $I_{R3} = 0.665\text{ mA}$, $I_{R4} = 0.89\text{ mA}$.
    \item Khi $V = 12\text{ V}$: $V_1 = 8.92\text{ V}$, $V_2 = 8.22\text{ V}$, $I_{R1} = 1.54\text{ mA}$, $I_{R2} = 1.115\text{ mA}$, $I_{R3} = 0.945\text{ mA}$, $I_{R4} = 1.37\text{ mA}$.
  \end{itemize}
\end{dapsobox}

\subsection{Mô phỏng PSpice}
\begin{itemize}[leftmargin=1.5em]
  \item Dùng hồ sơ \textbf{Bias Point}; lần lượt đặt $V = 8\text{ V}$ và $V = 12\text{ V}$, chụp màn hình có hiển thị dòng qua các điện trở và điện thế các nút $V_1$, $V_2$.
\end{itemize}
\simresult
\simfig{ex3_sim1}{Kết quả mô phỏng Bias Point của Bài tập 3 với $V = 8\text{ V}$}
\simfig{ex3_sim2}{Kết quả mô phỏng Bias Point của Bài tập 3 với $V = 12\text{ V}$}

\subsection{So sánh và nhận xét}
\begin{table}[H]
	\centering
	\caption{So sánh kết quả tính toán lý thuyết và mô phỏng PSpice cho Bài tập 3}
	\begin{tabular*}{\textwidth}{@{\extracolsep{\fill}}lcccc}
		\toprule
		\textbf{Đại lượng} & \multicolumn{2}{c}{\textbf{Lý thuyết (Theory)}} & \multicolumn{2}{c}{\textbf{Mô phỏng (PSpice)}} \\
		\cmidrule(lr){2-3} \cmidrule(lr){4-5}
		& $V = 8\text{ V}$ & $V = 12\text{ V}$ & $V = 8\text{ V}$ & $V = 12\text{ V}$ \\
		\midrule
		$I_{R1}\text{ (mA)}$ & $0.980$ & $1.540$ & $0.996$ & $1.553$ \\
		$I_{R2}\text{ (mA)}$ & $0.755$ & $1.115$ & $0.751$ & $1.112$ \\
		$I_{R3}\text{ (mA)}$ & $0.665$ & $0.945$ & $0.653$ & $0.935$ \\
		$I_{R4}\text{ (mA)}$ & $0.890$ & $1.370$ & $0.898$ & $1.377$ \\
		$V_1\text{ (V)}$     & $6.040$ & $8.920$ & $6.008$ & $8.894$ \\
		$V_2\text{ (V)}$     & $5.340$ & $8.220$ & $5.389$ & $8.260$ \\
		\bottomrule
	\end{tabular*}
\end{table}
